\documentclass[10pt,a4paper]{amsart}
\usepackage[margin=19mm]{geometry}
\usepackage{amsmath,amssymb,mathtools}
\usepackage[scr=rsfs,cal=euler]{mathalfa}
\usepackage{xfrac}
\usepackage[unicode,hidelinks,bookmarks=false]{hyperref}
\newcommand{\ie}{i.e.\ }
\newcommand{\cf}{cf.\ }
\newcommand{\resp}{respectively\ }
\newcommand{\Subsection}{\subsection}
\providecommand{\nobreakdash}{\nobreak\hbox{-}}
\setlength{\emergencystretch}{2em}
\multlinegap0pt
\usepackage{bm}
\usepackage{bbm}
\usepackage{dsfont}
\usepackage{xspace}
\usepackage{cancel}
\usepackage{mathtools}
\usepackage{amsthm}
\makeatletter
\@ifundefined{theorem}{\newtheorem{theorem}{Theorem}}{}
\@ifundefined{lemma}{\newtheorem{lemma}[theorem]{Lemma}}{}
\makeatother
\newcommand{\ldownarrow}{\Big\downarrow}
\title[Two combinatorial puzzles arising from the theory of Kohnert]{Two combinatorial puzzles arising from the theory of Kohnert polynomials}
\author{Theo Koss, Nicholas Mayers, Alex Moon}
\date{Source: arXiv:2509.17170v1 (CC BY 4.0)}
\begin{document}
\maketitle

\noindent\emph{Source and licence.} Two combinatorial puzzles arising from the theory of Kohnert polynomials. Version arXiv:2509.17170v1 (CC BY 4.0), distributed under the Creative Commons Attribution 4.0 International licence, as recorded in the arXiv licence metadata for this exact version and shown on the abstract page https://arxiv.org/abs/2509.17170v1. Full rights notice: licenses/arxiv-2509.17170-CC-BY-4.0.txt.\par\smallskip

\noindent\textit{Editorial scope.} Counted unit: the Lemma whose source label is \texttt{lem:minprophelp2} in the article (source0.tex lines 718--725, source label \texttt{lem:minprophelp2}), delivered together with the article's own complete original proof of it (source0.tex lines 726--751, one \texttt{proof} environment of 26 lines). No mathematics is written, repaired or re-derived: statement and proof are the article's own text line for line. Required context retained uncounted: lem:minprophelp (lines 608--614). Every definition, display and label of those lines is retained unchanged. Omitted from this package and not redistributed: 1--607, 615--717, 752--927. Nothing inside a retained excerpt is omitted or altered: every retained body is delivered exactly as the article's lines are written, so no omission receipt of any kind accompanies this package. Citations: the retained excerpts carry no citation command at all, so no bibliography entry of the article is transcribed and no reference is redistributed. Reference adaptation: none was needed and none was made; every reference inside the delivered unit resolves inside it, so no unresolved reference or citation is produced. Printed slips kept literal: no printed word, symbol, hypothesis or number of the counted statement or of its proof was altered. Layout and class: the article's own class is \texttt{article} with packages including geometry, fullpage, enumerate, authblk, hyperref, verbatim, section, amsthm, textcase, setspace, amssymb, amscd, lineno, amsmath; the wrapper is the lane's pinned \texttt{amsart} editorial wrapper at 10pt on A4, which restates the theorem-like environments the retained text needs and adds the article's own macro definitions that the retained text needs (\texttt{\textbackslash ldownarrow}), with extra wrapper packages (bm, bbm, dsfont, xspace, cancel, mathtools). The delivered numbering is the unit's own. Assets: no retained span contains a figure environment, a graphics-inclusion command, a PDF-page inclusion, a table or a TikZ picture; no image file is copied into this package, the package declares no asset dependency and no image file is redistributed with it. Licence: version arXiv:2509.17170v1 of this article is distributed under the Creative Commons Attribution 4.0 International licence, as recorded in the arXiv licence metadata for this exact version and shown on the abstract page https://arxiv.org/abs/2509.17170v1. A full rights notice is delivered at licenses/arxiv-2509.17170-CC-BY-4.0.txt. Characters: every retained excerpt is pure ASCII, so the delivered engine, XeLaTeX with its default Latin Modern text font, reports no missing character.
\begin{lemma}\label{lem:minprophelp}
    Let $D$ be a diagram with $\mathrm{cwt}(D)=(m_1,\hdots,m_n)$. Suppose that $M_j=\max\{m_i~|~j\le i\le n\}$ for $1\le j\le n$. Then for all $\tilde{D}\in\mathrm{Min}(D)$ and $1\le j\le n$,
    \begin{enumerate}
        \item[$(i)$] the cells in column $c$ of $\tilde{D}$ for $j\le c\le n$ occupy rows weakly below row $M_j$ and
        \item[$(ii)$] for each $1 \le \tilde{r} \le M_j$, there exists at least one column $c^*$ for which $j\le c^*\le n$ and $(\tilde{r},c^*)\in \tilde{D}$.
    \end{enumerate}
\end{lemma}

\begin{lemma}\label{lem:minprophelp2}
Let $D$ be a diagram and $c$ be a nonempty column of $D$. Then
\begin{enumerate}
    \item[$(i)$] for all $\tilde{c} < c$, $(r, \tilde{c}) \in \widehat{D}_c$ if and only if $(r, \tilde{c}) \in D$, i.e., $\widehat{D}_c$ and $D$ match in columns lying strictly to the left of column $c$;
    \item[$(ii)$] for all $\tilde{D}\preceq\widehat{D}_c$ and $\tilde{c}\ge c$, $(r,\tilde{c})\in \tilde{D}$ if and only if $(r,\tilde{c})\in\widehat{D}_c$, i.e., $\widehat{D}_c$ and $\tilde{D}$ match in columns lying weakly to the right of column $c$; and
    \item[$(iii)$] the topmost cell in column $c$ of $\widehat{D}_c$ occupies row $h(D,c)$.
\end{enumerate}
\end{lemma}

\begin{proof}
    Assume that $\mathrm{cwt}(D)=(m_1,\hdots,m_n)$. Moreover, for $i\in [n]$ with $m_i>0$, suppose that $(r_i,i)\in D$ is the topmost cell in column $i$ of $D$. Starting with $\widehat{D}_{n}$, note that by definition, if $r_n>m_n$, then 
    \begin{itemize}
        \item when applying Kohnert moves at rows $r_n$ down to $m_n+1$ in decreasing order the topmost cell in column $n$ of the corresponding diagram moves to a lower row; and
        \item after applying the Kohnert move at row $m_n+1$, the cells in column $n$ are bottom justified so that the remaining Kohnert moves applied in forming $\widehat{D}_{n}$ do nothing.
    \end{itemize}
    In the case that $r_n=m_n$, then the cells in column $n$ are bottom justified in $D$, $D=\widehat{D}_{n}$ and all Kohnert moves applied in forming $\widehat{D}_{n}$ do nothing. Consequently, in either case, $(i)$ through $(iii)$ hold.

    Now, assume the result holds for $\widehat{D}_{j}$ for $1<j\le n$. If $m_{j-1}=0$, then $\widehat{D}_{j-1}=\widehat{D}_{j}$ and, evidently, $(i)$ through $(iii)$ hold for $\widehat{D}_{j-1}$. So, assume that $m_{j-1}>0$. Let $M$ denote $\max\{m_i~|~j\le i\le n\}$. Note that since properties $(i)$-$(iii)$ hold in $\widehat{D}_{j}$, $$(\ast)\quad \text{the rightmost cells in rows $1\le r\le M$ of $\widehat{D}_j$ must contain no empty positions below}.$$ To see this, note that otherwise there would exist $(r,k)\in \widehat{D}_{j}$ with $j\le k\le n$ which is rightmost in its row and for which there is a max $r^*<r$ such that $(r^*,k)$ is an empty position in $\widehat{D}_{j}$; but then $$\widehat{D}_{j}\succeq D^*=\mathcal{K}(\widehat{D}_{j},r)=\widehat{D}_{j}\ldownarrow^{(r,k)}_{(r^*,k)},$$ contradicting our assumption that $\widehat{D}_{j}$ satisfies property $(ii)$. Now, recall that $\widehat{D}_{j-1}$ is formed from $\widehat{D}_{j}$ by sequentially applying Kohnert moves at rows $r_{j-1}$ down to $1$ in decreasing order. There are three cases.
    \bigskip

    \noindent
    \textbf{Case 1:} $M\ge r_{j-1}$. In this case, considering $(\ast)$, it follows that $\widehat{D}_{j}=\widehat{D}_{j-1}$ so that property $(i)$ holds for $\widehat{D}_{j-1}$ since it holds for $\widehat{D}_{j}$. Moreover, the topmost cell in column $j-1$ of $\widehat{D}_{j-1}$ occupies row $r_{j-1}=h(\widehat{D}_{j-1},j-1)$, so that property $(iii)$ holds for $\widehat{D}_{j-1}$. Finally, considering Lemma~\ref{lem:minprophelp} (i), since property $(ii)$ holds in $\widehat{D}_j$, it follows that all cells in columns $c$ of $\widehat{D}_j$ satisfying $j\le c\le n$ lie weakly below row $M$. Thus, since $r_{j-1}\le M$ and $\widehat{D}_{j}=\widehat{D}_{j-1}$, once again considering $(\ast)$, it follows that property $(ii)$ holds in $\widehat{D}_{j-1}$ as well.
    \bigskip

    \noindent
    \textbf{Case 2:} $r_{j-1}>M$. Assume that $r_{j-1}>m_{j-1}$; the other case following via a similar but simpler argument. In this case, while forming $\widehat{D}_{j-1}$ from $\widehat{D}_{j}$, when applying Kohnert moves at rows $r_{j-1}$ down to $m_{j-1}+1$ in decreasing order, the topmost cell in column $j-1$ of the corresponding diagram moves to a lower row. After applying the Kohnert move at row $m_{j-1}+1$, the cells of column $j-1$ are bottom justified and, considering $(\ast)$, the remaining Kohnert moves applied in forming $\widehat{D}_{j-1}$ do nothing. Consequently, property $(i)$ holds. Moreover, the topmost cell in column $j-1$ of $\widehat{D}_{j-1}$ occupies row $m_{j-1}=h(D,j-1)$ so that property $(iii)$ holds in $\widehat{D}_{j-1}$. Finally, note that the rightmost cell in rows $1$ up to $m_{j-1}>M$ in $\widehat{D}_{j-1}$ belong to either column $j-1$ which is bottom justified, or column $k$ for $j-1<k\le n$. Thus, once again considering $(\ast)$, $\mathcal{K}(\widehat{D}_{j-1},\tilde{r})=\widehat{D}_{j-1}$ for $1\le \tilde{r}\le m_{j-1}$. As, arguing in a similar manner to as in Case 1, all cells in columns $c$ for $j-1\le c\le n$ lie weakly below row $m_{j-1}$, it follows that $(ii)$ holds in $\widehat{D}_{j-1}$.
    \bigskip

    \noindent
    \textbf{Case 3:} $r_{j-1}>M>m_{j-1}$. In this case, while forming $\widehat{D}_{j-1}$ from $\widehat{D}_{j}$, when applying Kohnert moves at rows $r_{j-1}$ down to $M+1$ in decreasing order, the topmost cell in column $j-1$ of the corresponding diagram moves to a lower row. Since $M>m_{j-1}$, after applying the Kohnert move at row $M+1$, the topmost cell of column $j-1$ occupies row $M=h(D,j-1)$. Considering $(\ast)$, the remaining Kohnert moves applied in forming $\widehat{D}_{j-1}$ from $\widehat{D}_{j}$ do nothing. Consequently, we have that both properties $(i)$ and $(iii)$ hold in $\widehat{D}_{j-1}$. Now, arguing in a similar manner to as in Case 1, we have that all cells in columns $c$ for $j-1\le c\le n$ lie weakly below row $M$. Thus, since the rightmost cell in rows $1\le r\le M$ belong to a column $k$ for $j-1<k\le n$, once again considering $(\ast)$, it follows that property $(iii)$ holds in $\widehat{D}_{j-1}$.
    \bigskip

    \noindent
    The result follows.
\end{proof}

\end{document}
